% Generated draft: requires mathematical and source review.
\documentclass[11pt]{article}
\usepackage[margin=25mm]{geometry}
\usepackage{fontspec}
\setmainfont{Latin Modern Roman}
\usepackage{amsmath,amssymb,mathtools,unicode-math}
\setmathfont{Latin Modern Math}
\usepackage{enumitem,needspace,xurl,hyperref}
\setlength{\parindent}{0pt}
\setlength{\parskip}{5pt}
\setlength{\emergencystretch}{3em}
\newcommand{\N}{\mathbb{N}}
\newcommand{\Z}{\mathbb{Z}}
\newcommand{\Q}{\mathbb{Q}}
\newcommand{\R}{\mathbb{R}}
\newcommand{\C}{\mathbb{C}}
\newcommand{\F}{\mathbb{F}}
\newcommand{\A}{\mathbb{A}}
\newcommand{\PP}{\mathbb{P}}
\newcommand{\E}{\mathbb{E}}
\newcommand{\eps}{\varepsilon}
\newcommand{\dd}{\,\mathrm{d}}
\begin{document}
\section*{UnsolvedMath ID: 20000237}
\textbf{Source problem code:} catalogue entry
\subsection{Definitions and mathematical statement}
The original asks for a degeneration-based proof of the Mukhin--Tarasov--Varchenko theorem and an extension to cominuscule $G/P$. In the usual Grassmannian Schubert formulation, take $\operatorname{Gr}(k,n)$, distinct real points $t_i\in\PP^1(\R)$, and flags osculating a real rational normal curve. For Schubert conditions $\lambda^{(i)}$ with total codimension $k(n-k)$, the intended theorem asserts that the corresponding zero-dimensional intersection is transverse and all its points are real. The source does not specify which formulation or degeneration construction is intended, so this is an explicit conventional interpretation for review.

A cominuscule homogeneous space $G/P$ has the appropriate cominuscule parabolic structure, conventionally characterized by an abelian parabolic nilradical in the simple-group setting. The extension request must specify its osculating curves/flags and Schubert conditions; a universal all-real statement is only one possible reading, and is not supplied with hypotheses.

The research report gives a special case and an obstruction. It claims a two-cluster Pl\"ucker-coordinate degeneration proof for four Schubert divisors on $\operatorname{Gr}(2,4)$, with the binary-quartic discriminant expression
\[
 I=e_2^2-3e_1e_3+12e_4
\]
equal to a positive sum of three squares. Its coefficient/parameter convention for $e_i$ and the actual identity are not supplied. It then restricts the model to the principal-$\mathrm{SL}_2$ Lagrangian Grassmannian $\operatorname{LG}(2,4)=\operatorname{Sp}_4/P_2$, where it describes $I=0$ as the quadric equation.

For three distinct real osculating Schubert divisors in that setting, it reports exactly two transverse nonreal conjugate solutions: the residual-root quadratic has discriminant $-3$ times a squared Vandermonde. It attributes this to a known Sottile counterexample, refuting the naive unrestricted all-real cominuscule extension while explicitly leaving modified paired or type-dependent formulations untouched. Its candidate synthesis is that one quartic invariant exhibits the real Grassmannian degeneration and the nonreal Lagrangian behavior. This unreviewed special computation does not provide the requested general degeneration proof, nor classify every possible cominuscule extension.
\subsection{Short English statement}
Reprove the real Schubert theorem by degenerations and determine its appropriate cominuscule extensions. The report treats only a small Grassmannian case and a known Lagrangian counterexample to a naive all-real extension. Those results have narrower scope than the complete original request.
\subsection{Sources}
\begin{itemize}
\item Local UnsolvedMath registry, permanent ID $20000237$, source code \texttt{catalogue entry}; created and updated 2026-08-14. Supplied status \texttt{solved}; classification \texttt{SOLVED-BY-YOU}; difficulty L3; no proposer or year supplied. The classification is the registry's label, not a claim that this drafting session proved a theorem.
\item AIM Problem Lists, Algebraic geometry, workshop \emph{Degenerations in algebraic geometry}, Problem 10.6, source type PDF, \url{https://aimath.org/pastworkshops/degenalggeomproblems.pdf}. Canonical location \texttt{aim-algebraic-geometry-notes.json}, \texttt{notes[236]}, tag \texttt{problem}; statement status \texttt{exact}; original, clean, and displayed text agree. Extracted footer ``Notes by Tony Feng 6'' is retained as notes attribution/page information, not an extra hypothesis. Original remarks and literature field empty.
\item The original names Mukhin--Tarasov--Varchenko without an exact reference. The research summary attributes the Lagrangian counterexample to Sottile, without exact bibliography.
\item Supplied research attempt 1: result \texttt{counterexample}, run status \texttt{partially\_solved}, rigor label \texttt{proved}, low novelty confidence. Candidate contribution \texttt{invariant\_identity}: common binary-quartic mechanism for the $\operatorname{Gr}(2,4)$ degeneration and $\operatorname{LG}(2,4)$ obstruction. Explicit caveat: machine-generated research aid, not peer reviewed; independently verify before citation.
\item Dataset attribution: Ulam AI / UnsolvedMath Contributors, CC BY 4.0; original problem attribution remains to the supplied AIM workshop collection.
\end{itemize}
\end{document}
